% SOURCE docs/release/2026-10-29/source/docs/writing-repair/2026-10-29/src/4eb4d49e2544befe.tex CHAR 0
\chapter{What language models contribute to causal knowledge}
\label{la:frontier}

Section~\ref{scisyn:assistance} gives the principal recommendation for language
assistance. This chapter develops its evidential basis. It first distinguishes
knowledge, formal reasoning, mechanism proposals and workflow support, then
asks how published evaluations transfer to a banking decision. Exact examples
show how task mixtures, shared errors and unresolved causal models can change
that decision. The final comparisons examine consistency judges, orchestration,
literature extraction and the design of a versioned assistance experiment.

\section{Four different questions about cause and effect}

A banker asking why deposit balances fell can be asking for an explanation of
a recorded event, an estimate of a feasible rate intervention, a prediction of
next month's balance, or a diagnosis of an observation error. A fluent answer
can move between these objects without announcing the change. Evaluating assistance therefore starts by distinguishing the task: a proposed scientific question, evidence about a causal model,
a formal calculation conditional on that model, or a recommendation under a
specified loss. Success on one task does not establish success on the others.

The early frontier argument proposes using language models to organize domain
knowledge, draft graphs, critique assumptions and suggest robustness analyses.
It also records the ability to reproduce recognizable public data and
distinguishes transforming known information from discovering a new causal
structure. \citep{v12_la_2305_00050} This distinction is particularly consequential
in banking. A familiar statement about interest rates can be useful prior
knowledge, but the effect of a particular bank's deposit offer depends on its
population, competitors, eligibility rules and information clock. A remembered
relationship does not determine that intervention effect.

Formal benchmarks introduce a second object. CLadder constructs questions from
known causal structures with different verbalizations. Its results concern
answers to those generated questions, including historically specified models
and prompting methods. \citep{v12_la_2312_04350} Such a benchmark can reveal failures
to distinguish observation from intervention. It cannot establish that a model
elicits the right variables from an unfamiliar bank process, measures them
correctly, or returns calibrated uncertainty after fitting to dependent data.
These are additional stages of the proposed product.

Third, a language model can propose a causal mechanism for numerical analysis.
Model Discovery Agent illustrates a separation between executable mechanism
proposals and numerical fitting, evidence approximation and experiment design.
Its finite predictive evaluations explicitly do not establish causal
identification. \citep{v12_la_2608_09696} A mechanism may predict well because
several observationally equivalent models predict equally well. The bank must
still justify which intervention the model represents and which new experiment
would distinguish economically different mechanisms.

Finally, workflow assistance concerns selection and coordination of established
operations. Causal-learn+ describes assistance with diagnosis, assumptions,
method selection and interpretation while retaining algorithm outputs and
graph ambiguity. \citep{v12_la_2606_23608} This is a useful architectural direction.
It changes the analyst's interaction with numerical tools; it does not make a
method's premises true. The comparisons that follow keep these four contributions
separate while studying how they can work together.

\section{Text supplies a different information channel}

Write $D$ for the observed numerical records, $K$ for inspected documentary or
expert evidence, and $\mathcal M(D,K)$ for the models consistent with the
declared interpretation of both. For target $\psi$, a set-valued answer is
\[
 \Psi(D,K)=\{\psi(M):M\in\mathcal M(D,K)\}.
\]
This notation does not assert that finite data identify the population law or
that a statistical confidence procedure has been supplied. It describes which
models the analysis intends to consider. Statistical uncertainty about the law
must subsequently be propagated through the model set.

Additional reliable knowledge may exclude a model that the observed distribution
alone cannot exclude. A contract establishing that an invitation was randomized,
for example, has a different evidential role from a predictive correlation.
But a model-generated sentence repeating that randomization is common practice
does not establish what happened in the bank's historical process. The analysis
must identify the original source, the relevant population and dates, the
interpretation of its claim, and the person or measurement that supports it.

This also explains why retrieving more documents need not improve discovery.
LACR combines background knowledge, retrieved documents and optional numerical
information, and includes an UNKNOWN response. Its reported retrieval effects
differ across benchmarks, with adverse as well as favorable outcomes.
\citep{v12_la_2402_15301} Retrieval changes $K$; it does not certify $K$. Documents
may concern a different treatment definition, contain an association stated as
a cause, or duplicate the same original assertion. Ten repeated documents are
not ten independent confirmations of an edge.

An appropriate evidence object therefore retains the claim, source passage,
date, domain, units and proposed role. A passage may justify a variable
definition without justifying an arrow. A process document may establish that
one action occurs before another without excluding a common cause. A source
about another bank may suggest a hypothesis while leaving transport assumptions
unresolved. These roles determine the strength with which $K$ can restrict the
model class.

\section{Discovery methods use language judgments in different ways}

Three broad routes recur in the inspected literature: direct graph or relation
prediction, refinement of numerical discoveries, and priors or constraints used
within a search. \citep{v12_la_2402_11068} The distinction matters because it locates
where a wrong assertion enters the calculation.

In direct conditional-independence elicitation, chatPC obtains language-model
judgments from descriptions and context. Its statistical calculation concerns
the proportion of YES/NO model responses. \citep{v12_la_2406_07378} Repeated answers
can measure response variability under a prompt. Without a separate argument,
they are not repeated draws of the bank variables and do not yield a sampling
test of independence in the bank population. A highly stable wrong answer can
therefore be more dangerous than a visibly variable one.

Soft guidance provides a different failure boundary. The inspected CD-UCT
method combines language and mutual-information scores to guide graph search
without excluding otherwise valid DAGs. Language-only priors are not uniformly
beneficial in its reported comparisons. \citep{v12_la_2405_13551} Ancestral-path
constraints can likewise enter through a soft log-prior penalty, with reported
degradation under inaccurate prior information. \citep{v12_la_2306_16902} Softness
leaves room for the numerical criterion to resist the suggestion, but it does
not establish that the final graph probabilities are calibrated.

Hard restrictions can be scientifically appropriate when backed by a genuine
design fact. They are more consequential when inferred from prose: deleting a
candidate model deletes every effect value that only that model permits.
Statistical Causal Prompting converts model outputs into forced, forbidden or
unknown prior knowledge using method-specific mappings. The inspected examples
include failures as well as improvements. \citep{v12_la_2402_01454} The interface
must expose both the original assertion and the conversion rule, so a reviewer
can determine whether a numerical threshold has quietly become a scientific
exclusion.

Noisy language priors can also guide a sequential allocation of questions over
edges while retaining PAG representations. \citep{v12_la_2506_16234} This is closer
to a budgeted information problem than to asking for a finished graph in one
prompt. Its relevance to a bank still depends on whether the supplied knowledge
applies to the observed regime, whether latent structure is represented
correctly, and whether the retained ambiguity survives into the effect query.

\section{Temporal discovery is not a change of vocabulary}

A statement that collections activity affects arrears omits the action time,
available history, response horizon and possible feedback. Expanding a variable
name into a paragraph does not supply these missing indices. A temporal graph
needs explicit time labels and a rule governing contemporaneous edges, lag
limits, hidden histories, stationarity and interventions.

RealTCD combines a language-derived initial adjacency with numerical score
optimization over temporal graphs and intervention targets. Its setup includes
causal sufficiency, a known observational regime, finite lag and a DAG on
unrolled variables; its identification statement has further stated premises.
The inspected real-system assessment lacks a ground-truth DAG and evaluates
selected relationships through expert rules. \citep{v12_la_2404_14786} These
qualifications prevent a graph score from being read as a general guarantee
for a bank's latent, irregularly observed time series.

The Large Causal Models work is an adjacent numerical foundation-model route,
not simply a conversational language-model interface. The inspected model
assumes additive noise, stationarity, no latent confounders and no
contemporaneous effects, and does not detect or correct violations of these
assumptions. \citep{v12_la_2602_18662} Distinguishing the two routes avoids treating
the shared word ``model'' as a common scientific contract. One supplies semantic
proposals; the other learns a temporal mapping under a training distribution
and structural restrictions. Either can be useful within its declared regime.

For sequential servicing, a credible workflow starts with the decision history
and feasible intervention. It then asks whether the discovery output expresses
the required variables and whether a policy effect is identified over every
retained graph. An automatically generated lagged graph cannot resolve
unrecorded private information merely by assigning it a latent-node name.

\section{Interactive systems and agents add a review problem}

CausalChat combines graph suggestions, user editing and numerical updates.
Its small nonexpert study uses a fixed sequence of interfaces.
\citep{v12_la_2410_14146} Such evidence can motivate better interfaces, but it does
not isolate every learning, ordering or expertise effect. The proposed bank
evaluation needs qualified reviewers, an explicit task distribution and a
comparison that records the time and information supplied by people.

Debating agents provide another review structure. MAC uses affirmative,
negative and judge roles; its inspected comparisons also report a competitive
single-agent route for simpler graphs at lower cost. \citep{v12_la_2407_15073}
Several personas can share the same mistaken premise or model knowledge.
Agreement among them is not automatically independent corroboration. Their
usefulness should be measured through detected errors, retained correct
alternatives, decision consequences and total cost.

The oci-agent study instead emphasizes an accountable human principal over
actor/critic workflows and deterministic causal libraries. Its critic can
decline analyses with poor overlap. \citep{v12_la_2607_22443} The relevant outcome
then includes correct refusal and what remains possible for the requested
population. Reporting only completed analyses would hide precisely the
difficult cases in which an agent's judgment matters most.

An evidence-auditing agent faces a similar boundary. ARGUS audits reported
support for assumptions rather than adjudicating their truth. Its retrieval
and abstention tradeoffs and small human-label pilot leave substantial
generalization questions. \citep{v12_la_2609_30867} In banking, an omitted document,
a contradicted assumption and an assumption whose truth is unknown require
different responses. A documentation score should not become an identification
certificate.

\section{An evidence standard for the proposed banking product}

The synthesis supports a modular research program. Language assistance proposes
questions, variables, graph alternatives and evidence links. Formal operations
check declared structures and calculate identified functionals or sets in
supported regimes. Numerical procedures estimate those objects with their own
sampling and approximation arguments. Human and institutional evidence supply
facts about assignment, ownership and feasibility that neither interface can
manufacture.

Each interface needs an explicit refusal or unknown state. A retrieved passage
with the wrong unit should fail target review. A graph class outside a backend's
domain should return unsupported. A valid but uninformative effect set should
remain visible. A conflict between textual and numerical evidence should retain
both sources for investigation instead of selecting whichever yields a cleaner
answer.

CausalDS makes narrated mechanisms, observations and private scoring objects
separate parts of a benchmark, and its results distinguish symbolic answers,
uncertainty, abstention and cost. \citep{v12_la_2607_08093} The proposed evaluation
adopts that separation as a design principle. It does not borrow a published
model score as evidence for the local prototype. The remainder of this chapter
makes the comparison explicit: what must stay fixed for a published score to
support a banking capability, and what can fail when assistance stages are
combined? Later chapters develop translation, diagrams, discovery and workflow.
Empirical claims about a chosen backend still require a declared evaluation
on independent mechanisms.


% SOURCE docs/release/2026-10-29/source/docs/writing-repair/2026-10-29/src/0168fc9719536212.tex CHAR 0
\section{From a frontier result to a banking capability}
\label{exec13:frontier-synthesis}

The inspected literature gives reasons to invest in better assistance, but it
does not yield a single ranking of causal systems. A system that extracts
relations from documents, a system that answers a supplied structural-model
question and a system that chooses a treatment policy may all receive a score
called causal accuracy. Their inputs, permissible assumptions and correct
outputs differ. A product proposal must recover these differences before it
uses an experimental result to select software or allocate development effort.

For this purpose, describe an evaluation claim by
\begin{equation}
 \mathcal E=(\mathcal T,\mathcal I,\mathcal A,\mathcal O,\mathcal B,\mathcal L),
 \label{exec13:frontier-claim}
\end{equation}
where $\mathcal T$ is the task distribution, $\mathcal I$ the information
supplied to the system, $\mathcal A$ its scientific assumptions, $\mathcal O$
the oracle or adjudication rule, $\mathcal B$ the tool, repair and human budget,
and $\mathcal L$ the loss used to score its output. This is a proposed
comparison vocabulary, not a theorem that these quantities are observed without
error. It asks a practical question: which parts of the published experiment
would remain the same in the proposed banking use? A change in any component
requires an argument or a new experiment before the original score can carry
the same interpretation.

For example, a graph provided in the prompt belongs to $\mathcal I$. Correctly
calculating a back-door formula from it tests conditional reasoning. Inferring
that graph from account and contact records is a different task with different
assumptions. Likewise, a human who repairs a mistaken join changes the combined
workflow's information and cost. The repaired answer can be valuable, but
attributing its correctness entirely to the model would misdescribe the product
being evaluated. This motivates an explicit comparison of the strongest
available alternatives rather than an agent-only leaderboard.

\subsection{Favorable and adverse results locate different opportunities}

The earliest broad frontier analysis treats language models as a channel for
domain knowledge and a possible aid to graph construction, critique and
robustness analysis. It also probes memorization and prompt sensitivity. These
observations motivate assistance with familiar concepts while warning that a
recognizable public dataset can supply information absent from a fresh bank
problem. Neither renaming a variable nor recalling a benchmark answer
establishes a new identification result \citep{v12_la_2305_00050}.

CLadder gives a more controlled formal task: generated questions from specified
causal structures with alternative verbalizations. That design isolates
important distinctions among association, intervention and counterfactual
reasoning. The supplied structure and generator, however, are part of the
experimental information. A bank workflow must also determine whether its
variable, treatment and population definitions support that structure. The
paper's remaining language-to-formal-tool problem is therefore central to the
proposed product, not an incidental user-interface issue
\citep{v12_la_2312_04350}.

Two comparisons show why the proposal should retain nonagent and smaller-model
baselines. In CausalEval, adding function calls improves the reported GPT-4o
CLADDER point score from 69 to 76 relative to ordinary chain-of-thought, but
reduces its e-CARE score from 90 to 73. These are outcomes from the named
historical configurations and samples. They demonstrate variation across tasks,
not a general tool penalty or benefit. In the inspected relation-extraction
study, the DDI F1 values for Single-Prompt C and PubMedBERT are 67.4 and 87.9.
The task concerns binary pairwise relations and the fine-tuned system receives
labeled training examples. The result supports considering supervised
extraction where the document domain and annotation budget are suitable; it
does not compare whole causal pipelines at equal total cost
\citep{exec13_census_causaleval,exec13_census_graph_validation}.

Retrieval has similarly mixed consequences. LACR combines background knowledge,
documents and optional numerical information. Its ASIA skeleton F1 falls from 1
to approximately .73 when documents are added, while other reported cases
benefit. Documents can add an applicable mechanism, introduce a conflicting
definition or repeat the same unsupported assertion. The scientific question
is which information changed and why it was relevant, rather than how many
documents entered the prompt \citep{v12_la_2402_15301}.

Auditing assumptions is another promising but restricted task. ARGUS assesses
the reported support for difference-in-differences assumptions, uses injected
flaws as local ground truth and explicitly avoids adjudicating causal truth.
Its five-paper human-label pilot and in-sample calibration do not establish
the error rate on a new collection of studies. A bank can use such an assistant
to find missing discussion or contradictory passages while keeping assignment,
parallel trends and measurement facts open to scientific assessment
\citep{v12_la_2609_30867}.

CausalCollab makes the longitudinal issue concrete. It defines an incremental
stylistic effect in sequential human--model interactions under explicit
positivity and sequential exchangeability premises. Its intervention mixes
text versions through a specified functional construction. A proposed
collections-message rewriting policy is not automatically that intervention:
it must define which actual messages are delivered, to whom, and after which
history. Model-generated counterfactual edits or outcome labels in an
evaluation remain generated evidence. The theorem's premises, the estimator's
approximation and the operational rewriting policy must be reviewed separately
\citep{exec13_census_collaboration}.

These results favor a portfolio of bounded roles. Template-driven formalization,
supervised extraction, retrieved evidence, formal graph algorithms and numerical
causal estimators can each be a baseline or component. A language model is
valuable where it reduces a measured bottleneck while preserving the meaning
of the task. Its presence is not a requirement for every component. The
architecture chapter develops how to enforce those boundaries in handoffs;
the question here is what evidence can justify selecting them.

\subsection{Task composition can reverse a published ranking}

Let $s_{ak}$ be method $a$'s success probability on task stratum $k$ and $w_k$
the proportion of eligible requests in that stratum, with $\sum_k w_k=1$.
Under a common success definition, the overall probability is
\begin{equation}
 S_a(w)=\sum_k w_k s_{ak}.
 \label{exec13:frontier-mixture}
\end{equation}
Even when every $s_{ak}$ is known, ranking methods requires $w$. A pooled score
from one benchmark mixture is not a mixture-free property of a system. When
the strata use incompatible outcomes, they should first remain separate rather
than be forced into this formula.

As an arithmetic illustration, use the two reported CausalEval point-score
pairs just discussed. Give CLADDER weight $w$ and e-CARE weight $1-w$. The
FunctionCall-minus-chain-of-thought difference is
\begin{equation}
 \Delta(w)=.07w-.17(1-w)=.24w-.17.
 \label{exec13:frontier-rank}
\end{equation}
It changes sign at $w=17/24\simeq.7083$. At equal weights the difference is
$-.05$; at $w=.8$ it is $.022$. This reweighting uses published point scores
solely to illustrate composition. It supplies neither their sampling
uncertainty nor a bank's task mix. The original paired outcomes and
mechanism-level dependence would be needed for an uncertainty calculation
\citep{exec13_census_causaleval,exec13_frontier_exact}.

The same principle applies to missing tasks. If all evaluated requests assume
an observed confounder set, the experiment cannot establish performance on
latent-confounding requests, whatever their future bank frequency. If the
support and weights of the intended workload are unknown, report stratum-level
results and plausible planning mixtures. A later public or bank intake sample
can estimate the mix. Changing the weighting after observing which method wins
would instead change the decision rule being evaluated.

The task family should be the meaningful sampling unit. Paraphrases from the
same mechanism, multiple prompt seeds and repeated tool calls share information.
The repeated-call variance example in the evaluation chapter explains why
their count cannot simply replace the count of independent mechanisms. For a
matched comparison, pair methods on the same held-out mechanisms, preserve all
retries and report uncertainty at the level where the evaluation law supplies
independence. No variance estimate can be recovered from a published average
alone.

\subsection{Component correctness must compose through conditional events}

Suppose a request passes five stages: correct formulation $E_1$, observation
construction $E_2$, identification $E_3$, numerical execution $E_4$ and
reporting $E_5$. By the chain rule,
\begin{equation}
 P\!\left(\bigcap_{j=1}^5 E_j\right)
 =P(E_1)\prod_{j=2}^5P\!\left(E_j\mid\bigcap_{k<j}E_k\right).
 \label{exec13:frontier-chain}
\end{equation}
This identity does not require independent stages. It does require that each
factor refer to the same workflow and condition on all relevant preceding
successes. Marginal accuracies reported in five unrelated papers are not those
conditional probabilities. They may concern different data, expert information
and definitions of success. Multiplying them would manufacture an end-to-end
claim from incompatible experiments.

For an exact planning example, conditional probabilities
$(.96,.95,.90,.97,.98)$ yield $.78025248$. If these numbers are instead only
marginal success probabilities in a common request population, the joint
success is bounded by
\begin{equation}
 \max\!\left(0,1-\sum_{j=1}^5P(E_j^c)\right)
 \leq P\!\left(\bigcap_{j=1}^5 E_j\right)
 \leq\min_jP(E_j),
 \label{exec13:frontier-frechet}
\end{equation}
giving $.76$ to $.90$. The lower bound follows from the union bound for
failures; the upper follows because joint success implies every individual
success. Two authored populations of 100 requests attain the endpoints. In
one, failure sets of sizes 4, 5, 10, 3 and 2 are disjoint, leaving 76 complete
successes. In the other the failure sets are nested inside the ten largest
failures, leaving 90. Both have exactly the same stage marginals
\citep{exec13_frontier_exact}.

For the proposal, this means stage tests diagnose where to improve, while a
separate composition test measures the completed request. A syntactically
valid graph can propagate a wrong population through every later stage. A
correct estimator may precisely calculate the wrong target. Conversely,
several stages can fail on the same small subset, so adding their failure
counts overstates the number of affected requests. Record per-request
transitions and preserve the joint history needed to distinguish these cases.

\subsection{An answer key is another measured object}

Let $Z$ be a benchmark's recorded label, $T$ the scientifically intended answer
and $\widehat T$ the system output. For exact-match correctness,
\begin{equation}
 \left|1\{\widehat T=T\}-1\{\widehat T=Z\}\right|
 \leq1\{T\ne Z\}.
 \label{exec13:frontier-oracle}
\end{equation}
Whenever the label is correct, both correctness indicators are equal;
otherwise their difference is at most one. Taking expectations therefore
bounds the difference between scientific accuracy and label agreement by the
answer-key error probability. If a defensible upper bound is $\epsilon$,
reported agreement $a$ implies only
$\max(0,a-\epsilon)\leq P(\widehat T=T)\leq\min(1,a+\epsilon)$.
This statement assumes a meaningful common answer space; set-valued and
ambiguous tasks require an explicitly defined acceptance relation.

With agreement $.9$ and a stipulated label error $.1$, truth accuracy can be
$.8$ or 1 in two simple 100-case constructions. Errors in the key can either
conceal wrong outputs or penalize correct ones. This is an exact example, not
an allegation that any inspected benchmark has a ten-percent error rate
\citep{exec13_frontier_exact}. The practical lesson is to audit the generator,
label rule and edge cases before using a score as product evidence. An
incorrect IV-to-population-ATE oracle cannot be repaired by making an agent
match that oracle more reliably. Likewise, model-generated labels and
model-generated answers can share errors even when their agreement is high.

For public observations without known causal truth, use the information that
is actually available: variable definitions, temporal admissibility, numerical
reproduction, sensitivity to assumptions and agreement with a justified design.
Do not silently promote a published observational point estimate into ground
truth. For generated mechanisms, keep generator truth outside the analysis
interface and verify the reference calculation independently. An intentional
nonidentified case must accept appropriate refusal or a valid set, rather than
reward a precise but unjustified number.

\subsection{Causal uncertainty matters through the decision}

Return to a bank choosing an action $a\in\{0,1\}$. Let $V(a,M)$ be its net
value under causal model $M$. When the evidence retains class $\mathcal M$,
one way to display the decision consequence is regret
\begin{equation}
 R(a,M)=\max_{b\in\{0,1\}}V(b,M)-V(a,M).
 \label{exec13:frontier-regret}
\end{equation}
A model's graph accuracy and its decision regret can order systems differently.
An edge irrelevant to the intervention query may contribute graph loss without
changing the decision; a single uncertain edge on the treatment-outcome path
can reverse the preferred action. A credible evaluation reports both the
scientific object and the loss appropriate to the decision.

The architecture chapter gives two Gaussian worlds with an identical observed
distribution but effects of 0 and $1/2$. Rescale the effect to dollars so that
the two candidate gains from acting are 0 and 500, and set action cost to 250.
Not acting has value zero in both worlds. Acting has value $-250$ or 250.
Because the available information has the same law in both worlds, any
randomized rule acts with the same marginal probability $p$. Its expected
regrets are $250p$ and $250(1-p)$, so
\begin{equation}
 \inf_{0\leq p\leq1}\max\{250p,250(1-p)\}=125.
 \label{exec13:frontier-minimax}
\end{equation}
Equality holds at $p=1/2$; either deterministic rule has worst regret 250.
This restricted two-world calculation illustrates the cost of unresolved
information. It is not a recommendation to randomize real banking actions,
and the world class and loss are authored. It explains why a request for an
informative experiment can be more valuable than another confident narrative
\citep{exec13_architecture_exact,exec13_frontier_exact}.

The next experiment should target information capable of distinguishing models
with different preferred actions. A public simulator can check that the
proposed procedure chooses such an experiment under known truth. Public
observations can establish feasibility of a data pipeline or stress an
assumption. Neither proves that the same model class, costs or experiment
availability holds at a particular bank. Those transport questions belong to
a later study justified by the prototype's findings.

\subsection{A public-evidence portfolio with explicit transfer limits}

An affordable portfolio follows the four banking journeys already developed.
For funding, public panel and macro records support data-clock, aggregation
and reproducibility checks, while declared generated laws supply causal truth
for uncertainty tests. For servicing, a controlled history/action generator
tests positivity failures, feedback, policy definitions and the effect of
withheld state information. For discovery, known mechanisms should include
observationally equivalent graphs whose effects differ, so a confident single
graph is not always rewarded. For observation construction, public or generated
documents can test extraction and linkage while preserving the difference
between a faithfully extracted assertion and a true causal statement.

Each journey should compare a deterministic reference, an established
statistical route and the proposed assistance where its role is relevant.
Give all arms the same admissible information and account separately for
training labels, retrieval, tool use and human repair. Freeze mechanisms,
splits, budgets and scoring before final evaluation. A fixed held-out set can
be used for that registered comparison; repeated development against its
answers creates a new dependence and requires a new final evaluation design.
An unavailable or malformed answer remains in the assigned-request denominator.

The interpretation should stay attached to the actual evidence. A document
review can justify an integration candidate. A local exact example can expose
a failure mechanism. A controlled prototype study can compare methods under
its declared law. A live-model study adds evidence about a particular model
revision and task distribution. A bank study adds evidence about the target
process and its decision consequences. These are useful, different claims;
their separation allows the proposal to advance now without purchasing an
expensive experiment merely to satisfy a deployment gate.

The current research frontier therefore supports a concrete proposal: build
typed formulation and evidence assistance around an auditable causal core;
retain supervised, deterministic and formal alternatives; measure composition
and review costs; and direct new information acquisition toward consequential
ambiguity. It leaves open the reliability of language-mediated assumptions,
performance outside benchmark support, annotation validity and the economics
of sustained human review. Those are falsifiable next questions. A count of
papers, agents or completed reports does not answer them.
\label{exec13:frontier-synthesis-end}


% SOURCE docs/release/2026-10-29/source/docs/writing-repair/2026-10-29/src/d6cfedf5c0da0641.tex CHAR 0
\section{From a coherent answer to a defensible causal decision}
\label{exec17:frontier-start}

Three recent systems clarify different meanings of assistance. A verifier can
check whether an answer follows from the reasoning presented to it; an
orchestrator can choose and run a sequence of causal tools; a literature agent
can extract candidate relations and draw a diagram. These are valuable but
different contributions. A bank considering a retention offer needs an answer
to a further question: under which assumptions would that offer change a
specified outcome for a specified population? We compare the mechanisms of
Recursive Causal Audit, OrcheCause and EpiCausalX against this decision, keeping
reported demonstrations distinct from a proposed banking implementation.
\citep{exec17_frontier_chang,exec17_frontier_yang,exec17_frontier_liu}

\subsection{Process consistency and the scientific oracle}

Chang's Recursive Causal Audit, or RCA, addresses a specific failure: a model
can describe a problem in one way and then emit a contradictory final label,
particularly when challenged by a user. Its judge has no gold labels. It checks
the output schema, consistency among fields, whether the conclusion follows
from the structured reasoning, and whether a user hint dominates contrary
reasoning. Staged prompts demand more explicit variables, causal sketches and
assumptions when the earlier output fails. The paper explicitly recognizes
that a consistently incorrect argument can pass the judge. That limitation is
part of the mechanism's definition, rather than an accidental omission that
another formatting check can repair.~\citep{exec17_frontier_chang}

In a retention analysis, the agent might state that adviser discretion creates
confounding and then report the treated-minus-untreated retention difference
as an identified offer effect. A consistency judge could catch the conflict.
It would not necessarily catch a different response that consistently asserts
random assignment throughout, when advisers actually selected customers using
unrecorded risk. The latter response needs evidence about the assignment
process. Agreement between a rationale and a number is a necessary internal
check for a well-formed report, but does not supply that evidence.

The study uses a curated 454-vignette diagnostic suite and standardized pressure
protocols. Its limitations include subjective annotations, prompt dependence,
black-box explanations of model behavior and the absence of a component-level
ablation of RCA. The historical calls run through January 2026; they are not a
ranking of models available at the time of this proposal. There is also a
specific unresolved denominator issue: the statistical-analysis paragraph
describes L1 with 50 cases per class and 100 overall, while Table 7 totals 50 L1,
304 L2 and 100 L3 cases. We therefore do not reconstruct a pooled confidence
interval or repair the denominator by assumption.~\citep{exec17_frontier_chang}

The answer-transition metric requires equally careful scope. In the paper's
Table 4, the L2 labels are binary, VALID and FLAWED. For a fixed gold label,
changing an incorrect binary answer necessarily corrects it. Let $p$ denote
initial accuracy, $b$ the probability of becoming wrong conditional on an
initially correct answer, and $c$ the probability of becoming correct
conditional on an initially wrong answer. Then
\begin{equation}
 p_{\mathrm{final}}=p(1-b)+(1-p)c.
 \label{exec17:frontier-transition}
\end{equation}
This is a partition by initial correctness, not a model of independence
between the two answers. Substitution for three selected rounded rows in Table 4
agrees with their reported final accuracies to within 0.1 percentage point.
It is therefore unwarranted to call this particular binary table erroneous
merely because a more general answer vocabulary exists elsewhere in the
paper.~\citep{exec17_frontier_chang,exec17_frontier_exact}

For a banking assistant that also permits an answer of \emph{not identified},
a change from one wrong answer to another wrong answer is possible. To see
the difference, fix a task with true label $C$ and two wrong labels $W_1,W_2$.
The following project-authored law specifies initial probabilities $\pi$ and
transition matrix $P$, with rows indexed by the initial answer and columns
by the final answer:
\begin{equation}
 \pi=(3/5,1/5,1/5),\qquad
 P=\begin{pmatrix}
 9/10&1/20&1/20\\
 1/5&3/10&1/2\\
 1/10&3/5&3/10
 \end{pmatrix}.
 \label{exec17:frontier-multiclass}
\end{equation}
All rows sum to one. Among initially wrong answers, the probability of any
label change is $7/10$, whereas the probability of correction is only $3/20$.
The bad-flip probability is $1/10$. Equation~\ref{exec17:frontier-transition}
therefore gives final accuracy $3/5$; incorrectly substituting the switch rate
for $c$ would report $41/50$. Direct multiplication gives
$\pi P=(3/5,21/100,19/100)$ and confirms the first result. These are exact
constructed probabilities, not measured model performance.
\citep{exec17_frontier_exact}

The practical implication is to preserve a full answer transition table,
including refusal and tool failure, with a fixed target-specific gold label.
Whether refusal is correct depends on whether the question is identified
under the supplied assumptions. A pressure test should separately expose
unsupported endorsement, inappropriate rejection and appropriate withholding.
Its denominator must include the same assigned questions before and after
review; deleting failed or abstaining requests can change the apparent benefit.

\subsection{Confounding is not a proof of a zero effect}

A vignette in Chang's appendix describes treatment given to more severely ill
patients. It labels the positive causal assertion NO and describes the
association as spurious. The defensible lesson is that the supplied
association does not establish the asserted effect. Confounding by itself
does not establish that the treatment has no effect. A gold-label protocol
must distinguish an unjustified assertion from a false effect assertion.
\citep{exec17_frontier_chang}

Consider the following authored banking construction. Let $A$ indicate a
retention offer, $Y$ indicate retention over a fixed horizon, and $U$ be an
unrecorded customer characteristic. Draw $U$ as Bernoulli with probability
$1/2$. Independent exogenous uniform variables generate the offer and
outcome using the conditional probabilities
\begin{equation}
 \Pr(A=1\mid U=u)=\tfrac14+\tfrac12u,\qquad
 \Pr(Y=1\mid U=u,A=a)=\alpha+\beta u+\tau a.
 \label{exec17:frontier-scm}
\end{equation}
The outcome equation is retained under intervention; only the offer equation
is replaced. All conditional probabilities below lie strictly between zero
and one, so the comparison does not depend on deterministic treatment or
an empty observed treatment stratum.~\citep{exec17_frontier_exact}

\begin{center}
\begin{tabular}{lrrrrr}
\toprule
World & $\alpha$ & $\beta$ & $\tau$ & $E[Y(0)]$ & $E[Y(1)]$\\
\midrule
Positive effect & $7/20$ & $1/5$ & $1/10$ & $9/20$ & $11/20$\\
Zero effect & $3/10$ & $2/5$ & $0$ & $1/2$ & $1/2$\\
Negative effect & $1/4$ & $3/5$ & $-1/10$ & $11/20$ & $9/20$\\
\bottomrule
\end{tabular}
\end{center}

Each world has $\Pr(U=1\mid A=1)=3/4$ and
$\Pr(U=1\mid A=0)=1/4$. Consequently the observed conditional means are
\begin{equation}
 E[Y\mid A=0]=\alpha+\beta/4=2/5,\qquad
 E[Y\mid A=1]=\alpha+3\beta/4+\tau=3/5
 \label{exec17:frontier-observed}
\end{equation}
in all three worlds. The complete observed joint law of $(A,Y)$, not only
its two means, is identical: the cells $(0,0),(0,1),(1,0),(1,1)$ have
probabilities $3/10,1/5,1/5,3/10$. Under intervention,
$E[Y(a)]=\alpha+\beta/2+\tau a$, so the causal difference is $\tau$.
Thus the same positive observed difference of $1/5$ coexists with a
positive, zero or negative effect. Independent finite enumeration of the
threshold construction reproduces both the observed and intervention laws.
\citep{exec17_frontier_exact}

For a stipulated USD 100 value per additional retained customer and USD 5
cost per offer, incremental value is $100\tau-5$: USD 5, minus USD 5 or
minus USD 15 in the respective worlds. These assumed amounts illustrate why
an unjustified zero-effect conclusion and an unjustified positive-effect
conclusion can both mislead a decision. They are not estimates of bank prices
or returns. Even unlimited observations of $(A,Y)$ would leave these three
worlds observationally indistinguishable. The next investment must supply
defensible restrictions or new information, such as an appropriately designed
public or synthetic randomized demonstration before proposing an expensive
bank trial.~\citep{exec17_frontier_exact}

\subsection{Text constraints and the limits of graph-score tuning}

OrcheCause offers a complementary mechanism: task orchestration across causal
discovery, structural causal modeling, root-cause analysis and causal-effect
estimation. Its paper describes causal-learn, DoWhy, pgmpy and causalml as
underlying libraries. User text yields hard cause--effect constraints, while
external database text yields soft constraints. Preprocessing includes
imputation and removal of highly correlated features; hyperparameter search
can use BIC when a ground-truth graph is unavailable.
\citep{exec17_frontier_yang}

This distinction between hard and soft text makes a concrete design choice
visible. A hard requirement $A\to Y$ excludes every graph without that edge.
If the sentence actually describes a historical association, a mediated
effect, an intended policy or a different time scale, the true explanation
may already have been removed before a statistical test runs. A soft
constraint can permit disagreement with the data, but its weight is still
a modeling choice; an LLM's linguistic confidence is not automatically a
calibrated probability that the edge is true. In the proposed bank interface,
each relation therefore needs an explicit role: observed association,
mechanistic hypothesis, temporal prohibition, intervention definition or
accepted scientific restriction. The strength of a constraint must follow
its accepted meaning and scope.

The paper's table also argues against a single undifferentiated improvement
score. The following are the reported PC demonstrations; HPO denotes
hyperparameter optimization, DK domain knowledge from the original dataset
papers, and EK external knowledge. Lower structural Hamming distance (SHD)
and higher F1 are favorable under the paper's respective metrics.
\citep{exec17_frontier_yang}

\begin{center}
\begin{tabular}{llrrrr}
\toprule
Dataset & Metric & Baseline & HPO & DK+HPO & DK+EK+HPO\\
\midrule
ALARM & SHD &13&13&10&10\\
      & F1  &.927&.927&.940&.940\\
Sachs & SHD &21&23&14&10\\
      & F1  &.323&.343&.611&.688\\
\bottomrule
\end{tabular}
\end{center}

HPO alone makes Sachs SHD worse, from 21 to 23, while F1 improves. Thus
the paper's general statement of progressive improvement is too broad if
read as improvement of every metric at every stage. Nor does knowledge
drawn from the dataset's own original paper simulate a previously unknown
bank mechanism. The demonstration supports the value of investigating
knowledge-assisted discovery; it does not estimate deployment accuracy,
causal-effect coverage or the incremental value of each component on new
banking questions. Time-series analysis, high-dimensional low-sample
robustness and improved extraction from real documents are future directions
in this paper.~\citep{exec17_frontier_yang}

Observed fit cannot by itself solve this last step. The three complete laws
in Equation~\ref{exec17:frontier-observed} give the same likelihood for every
sample of the observed variables, while their intervention values differ.
No criterion computed solely from that observed distribution can select the
correct intervention value without additional restrictions. A graph score
may guide a search inside an assumed family; the product must still expose
the family and its exclusions. Similarly, removing a highly correlated
variable needs a causal-role review: predictive redundancy does not prove
that an adjustment variable or measurement proxy can be discarded.
\citep{exec17_frontier_exact}

For a prototype, the useful adoption is the modular orchestration idea and
the explicit distinction between text sources. Reproducing the reported
implementation requires a separate pinned-code and dependency review.
A first comparison would hold the dataset, target and compute budget fixed,
then compare no text, verified constraints, erroneous constraints and
contradictory documents. The outcome is not merely graph similarity: it
includes whether the effect is identified, whether a wrong constraint
changes the decision, and whether the system discloses the conflict.

\subsection{Literature extraction, diagrams and query-relative roles}

Liu and colleagues' EpiCausalX preprint describes a literature-driven
workflow with natural-language initiation and diagram construction. Its
methods extract relationship terms with regular expressions and infer
directions through heuristics informed by text and epidemiological theory.
The diagram tool checks acyclicity and uses topology to label exposures as
nodes without incoming edges and outcomes as nodes without outgoing edges.
The reported usability exercise includes 12 participants, split into two
groups of 6, and an average task initiation within two dialogue rounds.
These are workflow observations; they are not measurements of causal-effect
accuracy or coverage.~\citep{exec17_frontier_liu}

The topology rule illustrates an important boundary for business translation.
In the graph $U\to A$, $U\to Y$, $A\to Y$, the offer $A$ is still the
exposure of interest despite having an incoming edge. If retention $Y$
subsequently affects account value $V$, then $Y$ remains the outcome for a
retention question despite having an outgoing edge. Roles are relative to
the query and intervention, whereas root and sink status are graph
properties. The authored worlds above instantiate the first graph with
explicit positive probabilities. A drawing tool can validate that this
representation is acyclic; it cannot thereby prove the absence of omitted
common causes or identify the effect from observed $(A,Y)$ alone.
\citep{exec17_frontier_exact}

Literature extraction can nevertheless lower the cost of preparing a review.
A useful output is a candidate relation with its original passage, population,
time ordering, measurement definition and study design. An observational
association should remain an association until the proposed causal argument
and its assumptions have been checked. Negation, hypothetical wording,
direction reversal, selection and population mismatch deserve explicit
adversarial examples. The agent can propose a diagram and alternatives, then
pass a well-specified query to a separately reviewed identification routine.
The diagram is an argument to inspect.

The preprint itself places systematic expert review, real-world validation,
larger user studies and comparisons with manual work in its future program.
It also describes open-source release as a future step after initial
validation. Its displayed rights notice and the absence of a reviewed
software release prevent a direct-adoption decision here. The appropriate
proposal is to assess the documented workflow and develop independent
query-typing tests; the paper's usability result does not justify porting an
uninspected implementation or claiming current software availability.
\citep{exec17_frontier_liu}

\subsection{A public-evidence prototype that distinguishes these contributions}

The three mechanisms lead to a staged, falsifiable prototype. Its first
component is a collection of stipulated causal worlds and business questions,
including the equal-observation example, query-relative variable roles,
binary and three-label review tasks, and deliberately false textual
constraints. These fixtures have known answers because their laws are
specified. A second component uses public datasets and documented assumptions
to test tool invocation, provenance and the treatment of missing identification
information. Public observational data need not have known effects: the
expected response may be an identified set, sensitivity analysis or an
explicit inability to decide.

The comparison separates the three interventions. Adding a consistency judge
tests whether contradictory conclusions decrease at a declared review cost.
Adding text constraints tests whether correct information helps and wrong
information is exposed. Adding literature extraction tests whether supported
relations, populations and temporal definitions are preserved. Requests,
failed calls, refusals, changed assumptions and review time must remain in
the evaluation record. For generated mechanisms, replicate mechanisms rather
than treating repeated prompts about the same mechanism as independent
scientific evidence. Predeclared tolerances and held-out problem families
would be needed before estimating an improvement.

The resulting investment case would establish whether these separate
contributions merit a larger experiment. It would preserve the distinction
between coherent reasoning, a diagram, a successful tool call, a better
graph score and a justified banking decision. The attainable next step is
a narrowly specified prototype with known failures and honest uncertainty,
followed by bank-specific validation if that evidence warrants the cost.
\label{exec17:frontier-end}


% SOURCE docs/release/2026-10-29/source/docs/writing-repair/2026-10-29/src/7e04880baf364781.tex CHAR 0
\section{From a general LLM survey to a falsifiable banking assistance proposal}
\label{sec:exec24-llm-challenges}

A bank asking whether a language model can help design a collections intervention
needs an evaluation of a decision process, not merely a fluent answer. The model
might translate a manager's request, retrieve policy documents, propose variables
and arrows, generate analysis code, and explain a result. Each step can succeed
while the final intervention claim remains unjustified. Kaddour and colleagues'
\emph{Challenges and Applications of Large Language Models} supplies a historical
organizing framework: design choices, deployed behavior and scientific evaluation
create different limitations. The inspected original is the 72-page 19 July 2023
preprint, which explicitly calls itself an opinionated review of predominantly
text models. It supports that taxonomy and its authors' assessment, not an
exhaustive account of the October 2026 frontier, a causal algorithm, or a
reproduction of the many studies it cites.~\citep{exec24_llmchallenges_original}

The survey connects prompt sensitivity, hallucination, stale knowledge, brittle
evaluation, missing ablations and reproducibility. Its application discussion
distinguishes programming, reasoning, behavioral simulation and synthetic data.
Better code generation does not establish causal identification; a plausible
simulated borrower does not establish a bank population law. The following
construction converts these distinctions into explicit proposal criteria.
All numerical examples are authored exact examples, not observations of an LLM,
vendor quotes or estimates of bank value.~\citep{exec24_llmchallenges_reference}

\subsection{Translation must preserve the decision and the information clock}

Suppose a collections manager asks whether a reminder should be sent before a
payment deadline. A usable translation specifies the eligible accounts, the
decision time, the intervention actually available then, the outcome horizon
and the monetary loss function. A variable called \emph{customer engagement}
must resolve into a measurement, timestamp and missingness rule. Replies recorded
after the reminder cannot silently become pretreatment adjustment variables.
The proposed graph must state which arrows come from dated process evidence,
which reflect scientific assumptions, and which are uncertain alternatives.
A well-formed graph file cannot supply that missing evidence.

This is a proposed contract for the library, CLI and MCP interfaces. A translation
artifact should bind the request, source snapshots, decision time, variable
definitions, candidate graph and target query. A downstream identification request
should name that artifact and return assumptions and unresolved alternatives
with its formula or refusal. The explanation stage should cite the same artifact
identities. A new prompt, retrieved document or graph revision creates a new
candidate, requiring the appropriate checks again. Original model APIs, local
models or a deterministic questionnaire can provide the translation; no language
model is required to be the scientific authority.

The survey's prompt discussion reports sensitivity to wording, demonstrations
and their order. Its hallucination discussion distinguishes contradiction of a
source from an assertion not determined by that source. Unsupported, contradicted
and supported claims therefore need different statuses. An unsupported edge may
be a reasonable hypothesis, but repeated prompts cannot promote it into evidence.
Underlying method papers require their own inspected originals before their
technical results justify adoption.~\citep{exec24_llmchallenges_original}

\subsection{Repeated agreement does not remove shared errors}

Consider an authored error model for three agents reviewing the same proposed
arrow. Let $C$ denote a shared missing-context event, with $P(C)=q=1/10$.
Conditional on $C$, every agent is wrong. Otherwise, each agent is independently
wrong with probability $p=1/10$. This is a declared probability law, not a
measured agent ensemble. One agent's marginal error is
\[
 r=q+(1-q)p=\frac{19}{100}.
\]
The majority error under the declared law is
\[
 e_3=q+(1-q)\{3p^2(1-p)+p^3\}
     =\frac{313}{2500}=0.1252.
\]
Incorrectly treating the three marginal errors as independent gives
\[
 \widetilde e_3=3r^2(1-r)+r^3
     =\frac{47291}{500000}=0.094582.
\]
The difference is $15309/500000$, or 3.0618 percentage points. Majority voting
improves on the single-agent error here, but the independence calculation
overstates the improvement. Every ensemble still fails on $C$, so its error
cannot fall below $q$ under this shared-failure rule.
\citep{exec24_llmchallenges_reference}

For the bank, the shared event could omit the same post-treatment measurement
definition from every retrieved context. More agents do not obtain the omitted
timing evidence. The proposal separates agreement among generators from
corroboration by evidence and an independently specified semantic oracle.
An oracle may check a declared clock or known constructed causal world; it
cannot certify an unknown real-world arrow by inspecting strings. Diverse
prompts and providers are candidate controls whose benefit must be measured,
not assumed to imply independence. Human reviewers can share missing context.

\subsection{Generated people and generated laws answer different questions}

The survey reports both favorable and unfavorable results for models used as
behavioral participants, notes contamination and bias, and warns that a generated
distribution may not represent the real distribution. These are survey-level
findings. This proposal does not import an underlying experiment's percentage
or assert that its result generalizes to borrowers.
\citep{exec24_llmchallenges_original}

A separate authored example makes the decision consequence explicit. Assume two
borrower types, $Z=0,1$, with incremental gross benefits from a reminder of USD 20
and USD $-30$, respectively. Let the reminder cost USD 5 for either type. If the
share of type 1 in a target population is $w$, its average net benefit is
\[
 V(w)=(1-w)(20-5)+w(-30-5)=15-50w.
\]
A generated population with $w=1/5$ implies USD 5 per account, whereas another
target with $w=3/5$ implies USD $-15$. The break-even share is $w=3/10$.
These populations share the stipulated type-specific response law; changing only
their mixture reverses the decision. Without information about the target
mixture, even perfect evaluation on the generated population does not identify
target value. If only $w\in[0,1]$ is known, the compatible value interval is
$[-35,15]$ dollars per account.~\citep{exec24_llmchallenges_reference}

Public records and generated structural worlds remain useful before expensive
bank experiments. A public corpus tests retrieval, citation fidelity, variable
extraction and version handling under its own labels and rights. A constructed
structural model tests effect formulas, refusal on unidentified targets and
responses to deliberate clock violations against known answers. An LLM-generated
persona explores possible requests and failure cases. These uses have different
oracles. The last supplies neither the true borrower mixture nor truthful
counterfactual outcomes. Begin with these bounded tests and reserve transport,
behavioral fidelity and real monetary value for later evidence.

\subsection{Ablations must reveal interactions and count failed requests}

The survey criticizes missing controlled experiments and discusses
one-factor-at-a-time ablations as an improvement with cost limitations. An
authored factorial example shows why a banking prototype must also consider
interactions. Let two binary factors indicate retrieval and a separately
specified semantic check. Stipulate acceptance probabilities of $0.20$ with
neither, $0.50$ with retrieval alone, $0.40$ with the check alone and $0.90$ with
both. Retrieval's contrast is $0.30$ without the check and $0.50$ with it; the
check's contrast is $0.20$ without retrieval and $0.40$ with it. The interaction is
\[
 0.90-0.50-0.40+0.20=0.20.
\]
Adding isolated improvements to the baseline predicts $0.70$, missing the
stipulated combined $0.90$. Isolated comparisons do not identify the combined
result without an assumption about interaction.
\citep{exec24_llmchallenges_original,exec24_llmchallenges_reference}

An eligible prototype can compare the four configurations on the same frozen
tasks, retain paired task results and randomize execution order where it matters.
Its claim must name whether it concerns the finite suite or a sampled task
population. Caches, persistent agent memory and adaptive prompt repairs must be
isolated or modeled; otherwise earlier runs change later conditions. Split
evaluation by causal task family and relevant document or API version, not merely
random paraphrases of one scenario. Keep test prompts and oracle labels outside
development tuning.

Now stipulate 100 initiated requests: 20 fail operationally, 20 complete but
abstain, and 60 return numbers, of which 36 pass the declared numeric-output
audit. Conditional acceptance is $36/60=60\%$, while accepted numeric delivery
per initiated request is $36/100=36\%$. These answer different questions.
A justified abstention can itself be correct on an unidentified task, so the
prototype separately scores refusal correctness and unsafe numeric answers.
It must not call every abstention an error or remove operational failures from
the service denominator.

Stipulate USD 20 for model calls, USD 60 for review and USD 20 for infrastructure.
Cost per completed request is USD $100/80=1.25$, per numeric result USD $100/60$,
and per accepted numeric delivery USD $100/36=25/9$, approximately USD 2.78.
These include all initiated requests in their numerator and are illustrative
allocations, not current prices. Actual comparisons need the same task policy,
acceptance criteria, repair budget, refusal accounting and cost boundary.
\citep{exec24_llmchallenges_reference}

\subsection{Versioned assistance is a testable investment option}

The survey separates training repeatability from inference through a changing
API. It reports preliminary model-drift claims from other work and suggests
open models as a remedy. Those empirical numbers are not reproduced here.
Open weights permit useful control but do not alone fix data provenance,
hardware, precision, sampling or scientific validity. A run record needs the
model or provider revision, prompts, source snapshots, retrieval index, tool
schema, generation settings, oracle version, output and every repair attempt.
Where an API does not expose the needed revision, record that limit.
\citep{exec24_llmchallenges_original}

Candidate implementations include deterministic templates plus analyst review,
retrieval-assisted generation, external tools, or agents with bounded retries.
Fine-tuning adds data, rights and evaluation obligations; multi-agent review adds
dependence and cost obligations. Begin with the least complex baseline capable
of the declared task and admit additional machinery against measured conditional
benefit. A successful prototype supplies a reason to fund later testing in its
declared regime. It does not establish a supported integration, a complete
software survey, a frontier ranking or bank adoption. The survey strengthens
evaluation design while leaving native capability and bank value unestablished.

\subsection{What the comparison requires of the next question}

Taken together, the studies locate useful roles for assistance without assigning
one undifferentiated causal score. Retrieved text may change the admissible
model class; a judge may detect inconsistent reasoning; an orchestrator may
preserve a sequence of tool calls. Their evaluation must retain the target,
information supplied, scientific assumptions, oracle, budget and loss. A change
in task mixture, shared context or answer-key quality can alter the apparent
benefit even when each reported calculation is correct.

The resulting recommendation is to test bounded roles against deterministic,
supervised and formal alternatives on the same complete requests. Retain
conflicts, refusal and operational failure, and price the review work needed
to obtain a defensible result. Before any of those comparisons can be designed,
the banking question itself must be explicit. The next chapter develops that
translation, including population weights, intervention definitions and the
information available at the decision time.

